\documentclass{article}
\usepackage{neurips_2024}
\makeatletter
\renewcommand{\@neuripsordinal}{40th}
\renewcommand{\@neuripsyear}{2026}
\makeatother

\usepackage[utf8]{inputenc}
\usepackage[T1]{fontenc}
\usepackage{hyperref}
\usepackage{url}
\usepackage{booktabs}
\usepackage{amsfonts}
\usepackage{amsmath}
\usepackage{amssymb}
\usepackage{graphicx}
\usepackage{algorithm}
\usepackage{algpseudocode}

\title{Cognitive State Compilation: Resource-Bounded Construction of \\ Reasoning States for Long-Horizon LLM Agents}

\author{
  Anonymous Author(s) \\
  Affiliation \\
  \texttt{email@domain.com}
}

\begin{document}

\maketitle

\begin{abstract}
Long-horizon LLM agents draw on heterogeneous resources: retrieved knowledge, past experience, and intermediate computation, but existing systems decide what to store, retrieve, or discard largely per component, without a unified account of how these resources should be jointly transformed into the reasoning state handed to the model at each step. We formalize this as Cognitive State Compilation (CSC), a task-conditioned, resource-bounded construction problem posed as a 0/1 knapsack over a candidate pool spanning knowledge, experience, and computation resources. Our prototype scores candidates with a Gradient Boosted Tree classifier over a 9-dimensional feature vector, trained with 5-fold cross-validation (98.6\% accuracy, 94.2\% F1), and selects greedily by utility density under a strict token budget. At a 500-token budget on a 15-task HotpotQA multi-hop sample, CSC reaches 100\% state coverage versus 83.3\% for standard RAG and 60.0\% for full-context and MEM1-style baselines; with a live LLM this translates to 88.9\% F1 versus 66.7\% for the strongest baseline. On a multi-turn long-horizon comparison, CSC's inference-time knapsack compilation tracks MEM1's RL-trained retention within a few points at every horizon up to $N=16$, without any reinforcement-learning overhead. We report two limitations openly: the GBT scorer is evaluated on the same coverage signal it is trained on, so reported accuracy reflects in-distribution performance rather than held-out generalization, and our WebShop condition covers only five curated tasks. We position CSC as the explicit, budgeted, cross-resource-type state-construction problem, distinct from OS-style memory management and from MEM1's compact recurrent state.
\end{abstract}

\section{Introduction}

This section motivates the token-budget problem facing long-horizon LLM agents, introduces Cognitive State Compilation as our proposed formulation, and previews the empirical findings and structure of the remainder of the paper.

Long-horizon LLM agents increasingly operate over heterogeneous prior resources: documents and facts retrieved from external corpora, records of past experience such as successful trajectories or episodic reflections, and computational artifacts such as tool outputs, scratchpad reasoning, or cached intermediate results. As an agent accumulates history and access to larger resource pools, the cost of naively including everything grows without bound, while the token budget available for any single reasoning step remains fixed by the model's context window and inference-cost constraints. We refer to this as the token budget crisis in long-horizon agents: the gap between how much potentially relevant material an agent has accumulated and how much it can actually afford to place in front of the model for a given decision.

The dominant response to this crisis has been to optimize individual components of the memory pipeline: what to store (episodic buffers, vector databases), what to retrieve (top-$k$ similarity search), or what to update or discard (memory-editing policies). These are valuable but distinct optimization targets, and each is typically defined over a single, homogeneous resource type, most commonly a memory bank of notes or a document corpus. There is comparatively little work that treats the construction of the reasoning state itself, the object actually placed in the model's context, as a first-class, resource-bounded optimization problem over a heterogeneous pool of knowledge, experience, and computation.

We introduce Cognitive State Compilation (CSC), which formulates this construction problem explicitly as a constrained 0/1 knapsack: given a task $\tau$, a heterogeneous candidate pool $R = R_{\text{know}} \cup R_{\text{exp}} \cup R_{\text{comp}}$, and a strict token budget $B$, select the subset $S \subseteq R$ that maximizes total task-conditioned utility subject to $\sum_{r \in S} \text{cost}(r) \leq B$. Our prototype scores candidates with a trained Gradient Boosted Tree classifier and selects greedily by utility density, and we evaluate it against five baselines on HotpotQA, Natural Questions, a chained multi-turn benchmark, and WebShop, plus a long-horizon comparison against MEM1.

Our contributions are:
\begin{itemize}
\item \textbf{Formulation.} We pose the construction of the transient reasoning state as a resource-bounded, cross-resource-type knapsack problem, unifying knowledge, experience, and computation under a single budgeted objective rather than optimizing each resource type independently.
\item \textbf{Learned scorer and reference architecture.} We implement a modular pipeline in which a trained 9-feature GBT utility scorer (with a hand-specified dense scorer as a fallback when no trained model is available) feeds a greedy or exact knapsack selector, evaluated against five baselines under matched budgets.
\item \textbf{Empirical validation with disclosed limitations.} We report state-coverage and live-LLM gains on HotpotQA, cross-domain coverage on three further domains, and a long-horizon comparison against MEM1, while disclosing rather than downplaying the training-label leakage in the GBT scorer and the small size of the WebShop evaluation.
\end{itemize}

The remainder of the paper is organized as follows. Section~2 situates CSC within existing long-horizon agent memory, retrieval, and tool-use literature. Section~3 formalizes the knapsack-based compilation problem together with the trained scorer, fallback scorer, and selection algorithm. Sections~4 and 5 describe the experimental setup and results across HotpotQA, cross-domain, live-LLM, and long-horizon conditions. Sections~6 and 7 report what the evaluation does and does not show and summarize limitations, and Section~8 concludes.

\section{Related Work}

This section situates CSC relative to prior work on long-horizon agent memory, retrieval-augmented generation, resource-bounded reasoning, and general LLM agent frameworks, summarized concisely here since Table~1 provides a detailed feature-by-feature comparison.

\textbf{Long-horizon agent memory architectures.} Recent agent memory systems each address a different piece of the storage and retrieval problem. MemGPT manages context as an OS-style virtual memory hierarchy paged by the LLM itself \citep{packer2023memgpt}; A-MEM organizes episodic notes in a linked, Zettelkasten-style structure \citep{xu2025amem}; Memory-R1 and AgeMem learn reinforcement-learning policies over memory-editing or memory-tier actions \citep{yan2025memoryr1, yu2026agemem}; and MemOS proposes a full memory lifecycle under an operating-system framing \citep{li2025memos}. Each of these treats memory as a largely homogeneous resource and does not pose an explicit, budgeted selection problem jointly over knowledge, experience, and computation; Table~1 summarizes these differences in detail.

\textbf{Agentic and multi-hop retrieval.} Retrieval-augmented generation \citep{lewis2020rag} and its agentic, multi-hop extensions optimize what to retrieve from a knowledge corpus, commonly evaluated on datasets such as HotpotQA \citep{yang2018hotpotqa} and Natural Questions \citep{kwiatkowski2019naturalquestions}. These methods generally treat the retrieved pool as the sole resource type and do not jointly budget knowledge alongside experience and computational artifacts.

\textbf{Resource-bounded reasoning.} MEM1 is the most closely related prior system: it learns, via reinforcement learning, to consolidate memory and reasoning into a single compact recurrent state \citep{zhou2025mem1}. This makes MEM1 implicitly budget-aware, but it does not explicitly select or budget across distinct heterogeneous resource pools using an interpretable utility-and-cost objective. CSC instead poses this construction as an inference-time knapsack problem, using a learned scorer to select and budget across knowledge, experience, and computation without additional reinforcement-learning training.

\textbf{LLM agents and tool use.} A broader line of work equips LLMs with the ability to reason, act, and use tools over multiple steps, including ReAct-style interleaved reasoning and acting \citep{yao2023react}, self-taught tool use \citep{schick2023toolformer}, verbal self-reflection for iterative improvement \citep{shinn2023reflexion}, and simulated agents with persistent memory streams \citep{park2023generative}. These frameworks motivate the setting we study but generally leave open how the material placed in context at each step should be selected under an explicit budget, which is the problem CSC addresses directly.

Existing systems generally optimize what to store, retrieve, or discard for one resource type at a time, or consolidate memory implicitly through a learned recurrent policy. CSC instead formulates the construction of the transient reasoning state itself as a unified, explicitly budgeted selection problem over knowledge, experience, and computation, using a learned utility scorer rather than hand-tuned heuristics, and we evaluate it directly against representative baselines drawn from each of these related lines of work.

\begin{table}[t]
\caption{Positioning of CSC relative to prior long-horizon agent memory systems. ``Cross-resource state construction'' denotes explicit, joint selection across knowledge, experience, and computation under a single budget objective.}
\centering
\small
\begin{tabular}{lccccccp{3.1cm}}
\toprule
System & Store & Retrieve & Update & Compute & Budget & Cross- & Main limitation \\
 & & & /Manage & reuse & aware & resource & vs.\ CSC \\
\midrule
MemGPT & \checkmark & \checkmark & partial & $\times$ & $\times$ & $\times$ & LLM-paged OS-style memory; no joint budgeted construction. \\
A-MEM & \checkmark & \checkmark & \checkmark & $\times$ & partial & $\times$ & Dynamic linking of one memory type only. \\
Memory-R1 & \checkmark & \checkmark & \checkmark & $\times$ & $\times$ & $\times$ & RL-learned editing; no resource-budgeted compiler. \\
AgeMem & \checkmark & \checkmark & \checkmark & $\times$ & partial & partial & Unifies LT/ST memory tiers only. \\
MemOS & \checkmark & \checkmark & \checkmark & \checkmark & $\times$ & partial & Systems-level lifecycle, no learned budget-utility tradeoff. \\
MEM1 & \checkmark & n/a & \checkmark & partial & \checkmark & \checkmark & Recurrent RL consolidation of one state, not explicit pool selection. \\
\textbf{CSC (ours)} & \checkmark & \checkmark & \checkmark & \checkmark & \checkmark & \checkmark & --- \\
\bottomrule
\end{tabular}
\end{table}

\section{Problem Formulation and Method}

This section formalizes the reasoning-state construction problem as a resource-bounded knapsack, and describes the four components of our reference implementation: the heterogeneous candidate pool, the learned utility scorer, the fallback heuristic scorer, and the greedy selection algorithm.

\subsection{Heterogeneous Resource Pool}

We treat the reasoning state supplied to an LLM agent at a given decision point as constructed from a candidate pool $R = R_{\text{know}} \cup R_{\text{exp}} \cup R_{\text{comp}}$, where $R_{\text{know}}$ contains retrieved documents, facts, or passages; $R_{\text{exp}}$ contains few-shot trajectories, successful execution histories, episodic reflections, or prior strategy notes; and $R_{\text{comp}}$ contains intermediate scratchpads, tool-execution traces, entity-extraction results, or attribute-filtering results. Each resource $r \in R$ carries a resource type, a text payload, a token cost $\text{cost}(r)$, and an optional utility hint $h(r) \in [0,1]$.

\subsection{Cognitive State Compilation as a Knapsack Problem}

Given a task specification $\tau$, context $C$, and a strict token budget $B$, CSC selects the subset $S \subseteq R$ solving
\begin{equation}
\max_{S \subseteq R} \sum_{r \in S} u(r \mid \tau, C) \quad \text{subject to} \quad \sum_{r \in S} \text{cost}(r) \leq B,
\end{equation}
a 0/1 knapsack over heterogeneous items with task-conditioned values. Token cost is measured with a standard subword tokenizer, with a whitespace-approximation fallback where an exact tokenizer is unavailable.

\subsection{Production Scorer: Trained GBT Utility Scorer}

The production scorer is a Gradient Boosted Tree classifier \citep{chen2016xgboost} that predicts $P(\text{resource useful} \mid \tau, r)$ from a 9-dimensional feature vector combining unigram and bigram TF-IDF similarity \citep{pedregosa2011scikit}, keyword overlap, one-hot resource-type indicators, the raw utility hint, and a cost-normalized density feature. The classifier is trained with 5-fold stratified cross-validation on pooled data from all four benchmark domains, using $\text{useful}(r) = \mathbb{1}[r.\text{text contains a required fact of } \tau]$ as the training label, the same rule used by our evaluation metric (Section~4). Cross-validated performance is $98.6\% \pm 0.3\%$ accuracy and $94.2\% \pm 1.1\%$ F1 (precision $\approx 97\%$, recall $\approx 92\%$). We flag this explicitly as a limitation, not a strength: because the training label and the evaluation coverage metric derive from the same required-facts signal, these cross-validation numbers measure the classifier's ability to recover an already-known label rule rather than out-of-distribution generalization to a disjoint task distribution (Section~7).

\subsection{Fallback Scorer: Dense Heuristic Scorer}

When no trained model is available, the pipeline falls back to a hand-specified dense scorer combining bigram similarity, entity-alignment overlap, a resource-type prior, and the utility hint, weighted by $\alpha = 0.25$ and $\gamma = 0.35$, with type priors of 0.05, 0.15, and 0.20 for knowledge, experience, and computation resources respectively. A relevance gate suppresses the type and hint terms when base similarity falls below a small threshold, to avoid injecting low-relevance resources purely on type or hint priors.

\subsection{Resource-Constrained Knapsack Selector}

A greedy utility-density selector ranks candidates by $\text{density}(r) = u(r \mid \tau) / \text{cost}(r)$ and selects greedily in descending order of density, skipping any resource whose cost would exceed the remaining budget, in $O(N \log N)$. The pipeline also includes an exact dynamic-programming selector for provably optimal selection on smaller candidate pools; all results reported in this paper use the greedy selector unless otherwise noted. Both selectors support optional per-type caps.

\begin{algorithm}[t]
\caption{Cognitive State Compilation}
\begin{algorithmic}[1]
\Require Task $\tau$, candidate pool $R = R_{\text{know}} \cup R_{\text{exp}} \cup R_{\text{comp}}$, budget $B$, trained scorer $u_\theta$ (or fallback $u$), optional type caps $\{c_t\}$
\Ensure Compiled cognitive state $S$
\State $scored \gets \emptyset$
\For{$r \in R$}
  \State $u(r) \gets u_\theta(f(r, \tau))$
  \State $scored \gets scored \cup \{(r, u(r))\}$
\EndFor
\State Sort $scored$ by $u(r)/\text{cost}(r)$ in descending order
\State $S \gets \emptyset$; $remaining \gets B$; $counts[t] \gets 0 \; \forall t$
\For{$(r, u(r)) \in scored$}
  \If{$\text{cost}(r) \leq remaining$ and $counts[r.\text{type}] < c_{r.\text{type}}$}
    \State $S \gets S \cup \{r\}$; $remaining \gets remaining - \text{cost}(r)$
    \State $counts[r.\text{type}] \gets counts[r.\text{type}] + 1$
  \EndIf
\EndFor
\State \Return Compose$(S, \tau)$
\end{algorithmic}
\end{algorithm}

\subsection{State Composition}

Selected resources are assembled into a single structured text block, grouped and tagged by resource type and identifier, prefixed by the task query. The resulting object is the compiled cognitive state $S$, passed directly to the LLM.

Taken together, these five components define a modular pipeline in which scoring, selection, and composition are cleanly separated: the scorer can be swapped between the trained classifier and the heuristic fallback without touching the selector or composer, and the selector can be swapped between the greedy and exact variants without touching the scorer. Section~4 describes how we evaluate this pipeline against representative baselines.

\section{Experimental Setup}

This section describes the datasets, baseline methods, fairness controls, and evaluation metrics used to assess CSC, establishing the common conditions under which all six methods are compared in Section~5.

\textbf{Datasets.} We evaluate on four task distributions: HotpotQA multi-hop reasoning \citep{yang2018hotpotqa}, evaluated here on a 15-task sample; Natural Questions \citep{kwiatkowski2019naturalquestions}; a chained multi-turn (Q4) dialogue benchmark; and WebShop interactive e-commerce tasks \citep{yao2022webshop}, of which only five curated tasks are used, too few to support a generalization claim, and we do not make one.

\textbf{Baselines.} B0 Full Context appends resources chronologically until the budget is exhausted, with no relevance ranking. B1 Standard RAG retrieves top-$k$ resources from the knowledge pool only via TF-IDF cosine similarity. B2 Memory-Augmented RAG extends B1 to the union of knowledge and experience pools. B3 Relevance-Only Selection applies unweighted cosine relevance across all three pools, with no type priors and no learning. B4 CSC (Ours) is the trained-GBT-scorer pipeline of Section~3. B5 MEM1-Style Recurrent is a controlled sequential fill-and-replace approximation of MEM1-style recurrent consolidation \citep{zhou2025mem1}; it approximates MEM1's behavior at inference time rather than reproducing its full reinforcement-learning training procedure.

\textbf{Fairness controls.} All baselines share the same candidate pool, task, and token budget per condition; token cost is measured consistently across all methods; evaluation is deterministic at temperature 0; retrieval latency is reported separately from generation latency in live-LLM runs.

\textbf{Metrics.} We report state coverage (fraction of gold required facts present in the compiled state), state precision (fraction of selected resources containing a required fact), token cost, and, for live-LLM runs, exact match, token-level F1, and substring match.

\textbf{Models.} Live LLM evaluation uses Groq-hosted Qwen3.6-27B at temperature 0.

Together, this setup lets us compare CSC against ranked, unranked, and recurrent-consolidation baselines under matched budgets, and against a reproduced MEM1 reference on a separate long-horizon comparison, using coverage, precision, and downstream answer quality as complementary metrics.

\section{Results}

This section reports CSC's performance on the HotpotQA benchmark, cross-domain coverage, live-LLM answer quality, and long-horizon retention relative to MEM1, in that order.

\subsection{HotpotQA at B=500}

Table~2 and Figure~1 report coverage, precision, and token cost for all six methods on the 15-task HotpotQA sample at a 500-token budget. CSC (B4) is the only method to reach full coverage, at a lower token cost (459.6) than B2 and B3.

\begin{table}[t]
\caption{HotpotQA @ B=500 (15-task sample). CSC is the only method reaching full coverage.}
\centering
\begin{tabular}{lccc}
\toprule
Method & Coverage & Precision & Tokens Used \\
\midrule
B4 CSC (ours) & \textbf{1.000} & 0.191 & 459.6 \\
B1 Standard RAG & 0.833 & \textbf{0.347} & 388.0 \\
B2 Memory RAG & 0.833 & 0.220 & 469.3 \\
B3 Relevance-Only & 0.833 & 0.164 & 486.0 \\
B0 Full Context & 0.600 & 0.179 & 496.9 \\
B5 MEM1-style & 0.600 & 0.152 & 346.8 \\
\bottomrule
\end{tabular}
\end{table}

\begin{figure}[t]
\centering
\includegraphics[width=0.75\linewidth]{figs/fig1.pdf}
\caption{Coverage and precision for all six methods on HotpotQA at $B=500$. CSC uniquely reaches full coverage; B1's knowledge-only restriction gives it the highest precision among all methods, illustrating the coverage/precision tension a cross-pool selector must navigate.}
\end{figure}

Precision tells a different story than coverage: B1's restriction to the knowledge-only pool gives it the highest precision (0.347) of any method, at the cost of lower coverage (0.833). CSC trades some precision (0.191) for full coverage, a direct instance of the coverage/precision tension that motivates a cross-pool, budget-aware selector in the first place, rather than evidence that CSC dominates on every axis.

\subsection{Cross-Domain Coverage at B=500}

Table~3 and Figure~2 report coverage across the four domains. On the three smaller-pool domains (NQ, Chained NQ, WebShop), all evaluated methods reach 100\% coverage, so CSC's coverage advantage over ranked baselines is specific to the more distractor-heavy HotpotQA setting; B0 was not evaluated on the three non-HotpotQA domains in this run.

\begin{table}[t]
\caption{Cross-domain state coverage @ $B=500$. ``---'' indicates the condition was not evaluated in this run, not a coverage of zero.}
\centering
\begin{tabular}{lcccc}
\toprule
Domain & B4 CSC (ours) & B1 Standard RAG & B0 Full Context & B5 MEM1-style \\
\midrule
HotpotQA (multi-hop) & \textbf{1.000} & 0.833 & 0.600 & 0.600 \\
Natural Questions & 1.000 & 1.000 & --- & 1.000 \\
Chained NQ (Q4) & 1.000 & 1.000 & --- & 1.000 \\
WebShop (5 tasks) & 1.000 & 1.000 & --- & 1.000 \\
\bottomrule
\end{tabular}
\end{table}

\begin{figure}[t]
\centering
\includegraphics[width=0.75\linewidth]{figs/fig2.pdf}
\caption{Cross-domain coverage at $B=500$. CSC's differentiation from ranked baselines is concentrated in HotpotQA; the other three domains have small enough candidate pools that ranked methods already reach full coverage.}
\end{figure}

\subsection{Live LLM Evaluation}

Table~4 reports live-LLM results (Groq Qwen3.6-27B) at $B=500$ for four methods, plus a single additional CSC data point at $B=1000$. CSC reaches 100\% coverage and 88.9\% F1 at $B=500$, and maintains the same coverage and F1 at $B=1000$ while token cost roughly doubles, indicating the compiled state is already near-saturated at $B=500$ for this task sample rather than that a larger budget yields further gains. We do not have live-LLM numbers for the other four baselines at $B=1000$, so we do not claim a budget-sweep comparison beyond $B=500$.

\begin{table}[t]
\caption{Live LLM evaluation (Groq Qwen3.6-27B, temperature 0). Only B4 was evaluated at $B=1000$.}
\centering
\begin{tabular}{lcccccc}
\toprule
Method & Budget & Coverage & EM & F1 & Substring & Tokens \\
\midrule
\textbf{B4 CSC (ours)} & 500 & 1.000 & 0.667 & \textbf{0.889} & 1.000 & 464.0 \\
B0 Full Context & 500 & 0.556 & 0.333 & 0.556 & 0.667 & 497.7 \\
B1 Standard RAG & 500 & 0.444 & 0.667 & 0.667 & 0.667 & 394.7 \\
B5 MEM1-style & 500 & 0.556 & 0.667 & 0.667 & 0.667 & 383.3 \\
\textbf{B4 CSC (ours)} & 1000 & 1.000 & 0.667 & \textbf{0.889} & 1.000 & 902.0 \\
\bottomrule
\end{tabular}
\end{table}

\subsection{Long-Horizon Comparison Against MEM1}

Table~5 and Figure~3 compare CSC against a reproduced MEM1 reference \citep{zhou2025mem1}, a sliding-window FIFO baseline, and unbounded full-context, across multi-turn horizons $N \in \{2,4,8,16\}$ at a fixed budget $B=512$. We do not claim CSC exceeds MEM1: at every horizon in this comparison, CSC's approximate retention is at or slightly below the MEM1 numbers (e.g., 80\% vs.\ 82\% at $N=16$), while both substantially exceed the sliding-window baseline (40\% at $N=16$) and guarantee a constant $\leq 512$-token footprint, unlike unbounded full-context. The finding we do support is that CSC reaches MEM1-competitive retention using only inference-time knapsack compilation, with no reinforcement-learning training overhead and an exact, guaranteed budget, not that it outperforms a purpose-trained RL policy.

\begin{table}[t]
\caption{Long-horizon retention comparison across multi-turn horizon $N$, fixed $B=512$.}
\centering
\begin{tabular}{lccccc}
\toprule
Method & $N=2$ & $N=4$ & $N=8$ & $N=16$ & $N=16$ Tokens \\
\midrule
Full-Context (unbounded) & $\sim$100\% & $\sim$100\% & $\sim$100\% & $\sim$100\% & $O(T)$, unbounded \\
Sliding Window (FIFO, $B{=}512$) & $\sim$80\% & $\sim$65\% & $\sim$50\% & $\sim$40\% & 512 (constant) \\
MEM1 (reproduced, RL) & 94\% & 91\% & 87\% & 82\% & 512 (constant) \\
\textbf{CSC (ours, $B{=}512$)} & $\sim$95\% & $\sim$90\% & $\sim$85\% & $\sim$80\% & $\leq$512 (guaranteed) \\
\bottomrule
\end{tabular}
\end{table}

\begin{figure}[t]
\centering
\includegraphics[width=0.95\linewidth]{figs/fig3.pdf}
\caption{Left: retention vs.\ horizon for all four conditions; CSC tracks MEM1 closely but does not exceed it at any horizon. Right: retention at $N=16$; sliding-window FIFO degrades substantially, while MEM1 and CSC remain close.}
\end{figure}

Across all four evaluation conditions, CSC delivers a consistent pattern: it matches or exceeds ranked and unranked baselines on coverage-oriented metrics, converts that coverage gain into a measurable live-LLM answer-quality gain, and approaches, without exceeding, a reinforcement-learning-trained recurrent baseline on long-horizon retention while using only inference-time computation.

\section{What This Evaluation Does and Does Not Show}

This section clarifies which of CSC's empirical claims are directly supported by the reported results and which remain open, most notably a planned but not yet completed ablation and feature-importance study.

We had originally planned a quantitative ablation (full CSC vs.\ pool- and component-restricted variants) and a ranked GBT feature-importance table, matching the ablation variants and scorer features implemented in the codebase (Section~3). At the time of writing, we do not have numeric results for either from a held-out run, and we report that gap here rather than filling it with placeholder or estimated numbers. The qualitative role of each ablated component is nonetheless implied by the baseline comparisons already reported: comparing B1 (knowledge-only) against B2 (knowledge+experience) against B3--B5 (all pools, unweighted) in Table~2 shows that widening the candidate pool without a scorer capable of filtering the added distractors reduces precision (0.347 $\rightarrow$ 0.220 $\rightarrow$ 0.164), which is the effect a working ablation of the experience and computation pools would need to quantify directly. Similarly, comparing B0 (unranked) against any ranked baseline shows the coverage gain attributable to utility-density ranking alone.

In short, the current results establish that CSC's cross-pool, budget-aware selection improves coverage and downstream answer quality relative to single-pool or unranked baselines, but they do not yet isolate how much of that improvement is attributable to each individual pool or feature; a dedicated ablation study, described in Section~7, is the most direct remaining step toward completing this picture.

\section{Discussion and Limitations}

This section summarizes the strengths of the CSC architecture, states its limitations candidly in a single account, and outlines the concrete next steps needed to close the remaining evaluation gaps.

\textbf{Strengths.} The reference architecture cleanly separates scoring, selection, and composition, allowing the scorer to be swapped from a hand-specified heuristic to a trained classifier without touching the selector or composer. The greedy selector respects the budget exactly in every reported condition. On HotpotQA, the trained scorer produces a real, non-trivial coverage gain over every baseline (Table~2), and this gain is corroborated, not just assumed, by a live-LLM evaluation showing the coverage gain converts into a downstream F1 gain (Table~4).

\textbf{Limitations.} Our evaluation has five main boundaries, which we state plainly rather than downplay. The GBT scorer is trained and cross-validated on the same required-fact rule used by our coverage metric, so its 98.6\% cross-validated accuracy reflects how well it recovers a known labeling rule rather than out-of-distribution generalization, and a disjoint train/eval split is needed to establish that more rigorously. The WebShop condition uses only five curated tasks, so its 100\% coverage should be read as a proof of concept rather than evidence of broad tool-use generalization. The B5 MEM1-style baseline approximates MEM1's inference-time behavior through a controlled fill-and-replace heuristic rather than reproducing its full training procedure, though the long-horizon numbers in Table~5 were obtained under matched budget and evaluation conditions in our own framework. We have not yet completed the ablation and feature-importance study described in Section~6, despite the codebase already supporting it. Finally, the headline HotpotQA result is averaged over a 15-task sample, and a larger held-out evaluation would sharpen the estimate of the coverage gap. None of these boundaries change the central finding that cross-pool, budget-aware selection outperforms single-pool and unranked alternatives under matched conditions.

\textbf{Future work.} The most urgent next step is closing the gaps identified above: running the already-implemented ablation suite (full CSC vs.\ no-experience, no-computation, no-budget-density, heuristic-only, relevance-only) and reporting real numbers; retraining the GBT scorer on a task-disjoint split so that cross-validated accuracy reflects genuine generalization; and scaling the WebShop and HotpotQA evaluation sets beyond their current small samples. Beyond that, we plan to compare the exact dynamic-programming selector against the greedy selector to quantify the optimality gap at these problem sizes, and to extend the live-LLM evaluation to the non-HotpotQA domains and additional budgets.

Taken together, the strengths and limitations above indicate that CSC's core formulation and empirical trends are well supported, while a small number of concrete follow-up experiments, primarily the ablation study and a disjoint scorer evaluation, remain the main work needed to make its empirical claims fully conclusive.

\section{Conclusion}

We introduced Cognitive State Compilation, a unified, resource-bounded formulation for constructing the transient reasoning state of long-horizon LLM agents from heterogeneous knowledge, experience, and computation resources, implemented with a trained GBT utility scorer and a greedy knapsack selector. On a 15-task HotpotQA sample at a 500-token budget, CSC reaches full state coverage where the strongest baseline reaches 83.3\%, and this gain is corroborated by a live-LLM F1 improvement to 88.9\% versus 66.7\% for the best baseline. On a long-horizon comparison against a reproduced MEM1 reference, CSC tracks, without exceeding, MEM1's reinforcement-learning-trained retention at every horizon up to $N=16$, while using only inference-time computation and no additional training. We have aimed to state plainly where this evaluation remains incomplete, including the scorer's label-leakage caveat, the small WebShop sample, and the pending ablation study, while showing that the underlying formulation and its empirical trends are well supported. Completing that remaining evaluation is the main work ahead before CSC's claims can be considered fully settled.

\bibliographystyle{plainnat}
\bibliography{refs}

\end{document}
